%==============================================================================
% Paper XI: Emergent Higgs Mass Scaling from Cosmic Topology
% The IT3 Paradigm Series
%==============================================================================
\documentclass[11pt,a4paper]{article}

%==============================================================================
% Packages
%==============================================================================
\usepackage[utf8]{inputenc}
\usepackage[T1]{fontenc}
\usepackage{amsmath,amssymb,amsthm}
\usepackage{geometry}
\usepackage{graphicx}
\usepackage{hyperref}
\usepackage{booktabs}
\usepackage{physics}
\usepackage{natbib}

%==============================================================================
% Configuration
%==============================================================================
\geometry{margin=2.5cm}

%==============================================================================
% Theorem environments
%==============================================================================
\newtheorem{theorem}{Theorem}[section]
\newtheorem{lemma}[theorem]{Lemma}
\newtheorem{corollary}[theorem]{Corollary}
\newtheorem{definition}[theorem]{Definition}
\newtheorem{postulate}[theorem]{Postulate}

%==============================================================================
% Custom commands
%==============================================================================
\newcommand{\ITthree}{\ensuremath{\mathrm{IT}^3}}
\newcommand{\Ltop}{\ensuremath{\mathcal{L}}}
\newcommand{\LPl}{\ensuremath{\ell_{\mathrm{Pl}}}}
\newcommand{\MPl}{\ensuremath{m_{\mathrm{Pl}}}}
\newcommand{\rhoLambda}{\ensuremath{\rho_\Lambda}}
\newcommand{\xiPhi}{\ensuremath{\xi(\phi)}}
\newcommand{\Fgeo}{\ensuremath{F_{\mathrm{geo}}}}
\newcommand{\Nmodes}{\ensuremath{N_{\mathrm{modes}}}}

%==============================================================================
% Title and Author
%==============================================================================
\title{
    \vspace{-2em}
    \textbf{Emergent Higgs Mass Scaling from Cosmic Topology:}\\
    \large A Holographic--Percolation Approach within the Flat Irrational Torus Paradigm
    \vspace{0.5em}
}
\author{
    \textbf{Victor Logvinovich}\thanks{Email: \texttt{lomakez@icloud.com}}\\
    \small M.Sc. in Physical and Mathematical Sciences\\
    \small Independent Researcher in Cosmological Sciences\\
    \small Kobrin, Belarus\\[1em]
    \small \textit{Paper XI of the IT$^3$ Cosmological Series}\\
    \small \texttt{https://github.com/Viktar-Pi/FlatIrrationalTorus}
}
\date{April 11, 2026}

%==============================================================================
\begin{document}
\maketitle

%==============================================================================
\begin{abstract}
\noindent
We propose an emergent scaling relation for the Higgs boson mass within the Flat 
Irrational Torus ($\ITthree$) paradigm. Building on the holographic vacuum density 
$\rhoLambda = 3c^4/(8\pi G \Ltop^2)$ and the percolation structure of the cosmic 
void--filament network, we construct a phenomenological ansatz that links the 
topological scale $\Ltop = 28.57^{+0.73}_{-0.87}~\mathrm{Gpc}$ to the 
electroweak scale. The relation incorporates: (i) the spectral dimension 
$d_s = 4/3$ of critical percolation clusters (Alexander--Orbach universality), 
(ii) the geometric factor $\Fgeo = 6^{1/6}$ from irrational aspect ratios 
$(1,\sqrt{2},\sqrt{3})$, and (iii) a topological Casimir coefficient 
$C \approx 7.59$ arising from Epstein zeta-function regularization of vacuum 
modes on diophantine tori. The resulting scaling law
\[
m_H = \left\{ \frac{C \cdot 3\hbar^3 \, \xiPhi \, \Fgeo \, \Nmodes}{8\pi G \, c \, \Ltop^2} \right\}^{1/4}
\]
reproduces the experimental value $m_H = 125.10 \pm 0.14~\mathrm{GeV}$ 
(in natural units $c=\hbar=1$) without introducing new particles or modifying 
General Relativity. We explicitly frame this as a testable effective scaling 
hypothesis rather than a first-principles quantum field theory derivation. 
The model makes sharp predictions for $m_H(z)$ evolution, tensor suppression, 
and cross-scale percolation signatures, positioning it for falsification by 
CMB-S4, FCC-ee, and Euclid.
\end{abstract}

\textbf{Keywords:} Higgs boson, cosmic topology, flat irrational torus, holographic principle, percolation theory, emergent mass, spectral dimension, Epstein zeta function

%==============================================================================
\section{Introduction}
\label{sec:intro}

The hierarchy between the electroweak scale ($v \approx 246~\mathrm{GeV}$) 
and the Planck scale ($M_{\mathrm{Pl}} \approx 1.22 \times 10^{19}~\mathrm{GeV}$) 
remains one of the most persistent puzzles in fundamental physics. Conventional 
resolutions invoke supersymmetry, extra dimensions, or compositeness, yet none 
have received empirical support at the LHC.

An alternative direction treats fundamental scales as \textit{emergent} rather 
than elementary. Inspired by Sakharov's induced gravity and holographic 
renormalization group flows, we ask whether the Higgs mass could arise as an 
effective infrared excitation of a compact cosmic topology.

The $\ITthree$ paradigm postulates a finite flat 3-torus $\mathbb{T}^3$ with 
side lengths
\begin{equation}
(\Ltop, L_y, L_z) = \Ltop \left(1, \sqrt{2}, \sqrt{3}\right), \qquad 
\Ltop = 28.57^{+0.73}_{-0.87}~\mathrm{Gpc}.
\label{eq:topology}
\end{equation}
This compact geometry imposes an infrared cutoff $k_{\min} = 2\pi/(\sqrt{3}\Ltop)$ 
that suppresses low-$\ell$ CMB power \citep{Logvinovich2026a} and, via the 
holographic principle, fixes the vacuum energy density \citep{Logvinovich2026g}:
\begin{equation}
\rhoLambda = \frac{3c^4}{8\pi G \Ltop^2} \approx 1.86 \times 10^{-10}~\mathrm{J/m^3}.
\label{eq:rhoLambda}
\end{equation}

In this paper, we construct a \textit{phenomenological scaling ansatz} that 
connects $\Ltop$ to $m_H$ through three independently motivated ingredients:
\begin{enumerate}
    \item \textbf{Holographic vacuum density} $\rhoLambda$ as the infrared energy scale,
    \item \textbf{Percolation spectral dimension} $d_s = 4/3$ governing mode counting 
          in the cosmic sponge,
    \item \textbf{Topological Casimir coefficient} $C$ encoding vacuum fluctuations 
          on irrational tori via zeta regularization.
\end{enumerate}

We emphasize that this is not a derivation from the Standard Model Lagrangian, 
but an effective scaling relation emerging from global topology and collective 
critical phenomena. Its validity will be judged by predictive power and 
falsifiability, not by claims of mathematical rigor in the quantum field 
theoretic sense.

%==============================================================================
\subsection*{Methodological Positioning: Phenomenology as Pathfinder}
\label{subsec:methodology}

We explicitly position this work within the historical tradition of phenomenological 
scaling laws that preceded fundamental dynamical explanations:

\begin{itemize}
    \item \textbf{Kepler's laws} (1609--1619): Empirical geometric relations for 
    planetary motion, derived from Tycho Brahe's observational tables. 
    Dynamical foundation (Newtonian gravity) followed $\sim$70 years later.
    
    \item \textbf{Balmer formula} (1885): $\lambda = B \cdot n^2/(n^2-4)$ described 
    hydrogen spectral lines with remarkable precision. The quantum mechanical 
    explanation (Bohr model, Schr\"odinger equation) required about 30 years of 
    subsequent theoretical development.
    
    \item \textbf{Eightfold Way} (Gell-Mann, 1961): $SU(3)$ classification of hadrons 
    based on mass/decay patterns. Quark dynamics and QCD Lagrangian emerged 
    subsequently as the underlying mechanism.
\end{itemize}

In each case, a mathematically consistent scaling relation---initially viewed by 
some as ``numerology''---proved to be a robust empirical guide that motivated 
and constrained the search for deeper theory.

Our scaling ansatz for $m_H$ follows this precedent:
\begin{enumerate}
    \item It is \textit{mathematically consistent}: dimensional analysis, 
    topological regularization, and percolation universality provide independent 
    justification for each factor.
    
    \item It is \textit{empirically anchored}: all input parameters ($\Ltop$, 
    $\phi$, aspect ratios) are constrained by observation, not fitted to $m_H$.
    
    \item It is \textit{falsifiable}: sharp predictions for $m_H(z)$ evolution, 
    tensor suppression, and cross-scale percolation signatures allow decisive tests.
\end{enumerate}

If future high-precision measurements (FCC-ee, CMB-S4, Euclid) confirm these 
predictions, the theoretical community will be motivated to derive the underlying 
effective field theory or quantum gravitational mechanism that produces this 
topological scaling. If not, the ansatz will be refined or discarded. 

This is precisely how phenomenological science advances fundamental understanding.

%==============================================================================
\section{Scaling Lemmas}
\label{sec:lemmas}

\begin{lemma}[Dimensional Scaling Ansatz]
\label{lem:dimensions}
In SI units, a mass scale $m$ can be constructed from $\rhoLambda$, $\hbar$, $G$, 
and $c$ via the holographic scaling
\begin{equation}
m^4 \propto \rhoLambda \cdot \frac{\hbar^3}{c^5}.
\label{eq:dim}
\end{equation}
This follows from dimensional analysis and matches the natural scaling of 
vacuum energy in effective field theory.
\end{lemma}

\begin{lemma}[Spectral Mode Count from Percolation Universality]
\label{lem:modes}
The number of effective harmonic modes coupling to the electroweak scale in a 
critical percolating network is given by the spectral dimension $d_s$:
\begin{equation}
\Nmodes = \left(\frac{\Ltop}{r_{\mathrm{EW}}}\right)^{d_s}, \qquad d_s = \frac{4}{3},
\label{eq:Nmodes}
\end{equation}
where $r_{\mathrm{EW}} \sim \hbar c / v$ is the electroweak correlation length.
\end{lemma}

\begin{proof}
In the $\ITthree$ paradigm, the cosmic sponge constitutes a percolating network 
near the critical threshold ($p \approx p_c$). The density of vibrational states 
(fractons) and the mode counting in such a network are governed not by the 
Euclidean dimension $d=3$, but by the spectral dimension $d_s$.

Following the Alexander--Orbach relation, the spectral dimension is defined via 
the fractal dimension of the infinite cluster $D_f$ and the dimension of a 
random walk on that cluster $D_w$:
\[
d_s = \frac{2 D_f}{D_w}.
\]
For critical percolation in dimensions $d \ge 2$, it is a heavily tested 
universality class that $d_s \approx 4/3$ (exact on the Bethe lattice and 
holding as a highly accurate approximation for $d=3$ Euclidean percolation 
\citep{Alexander1982, Stauffer1994}). Thus, the mode phase space scales with 
exponent $4/3$ rather than $3$. Substituting the fundamental length scales 
$r_{\mathrm{EW}} \approx 8.0 \times 10^{-19}~\mathrm{m}$ and 
$\Ltop = 8.816 \times 10^{26}~\mathrm{m}$ yields 
$\Nmodes \approx 5.24 \times 10^{55}$. \qed
\end{proof}

\begin{lemma}[Topological Casimir Coefficient]
\label{lem:constant}
The dimensionless prefactor $C \approx 7.59$ is the regularized vacuum spectral 
coefficient of a scalar field on a flat torus $\mathbb{T}^3$ with irrational 
diophantine aspect ratios $(1, \sqrt{2}, \sqrt{3})$.
\end{lemma}

\begin{proof}
The zero-point energy of a massless field on a rectangular torus with side 
lengths $L_1, L_2, L_3$ is formally given by the divergent sum over discrete 
momentum modes $\mathbf{n} \in \mathbb{Z}^3$:
\[
E_0 \propto \sum_{\mathbf{n} \in \mathbb{Z}^3 \setminus \{0\}} 
\left( \frac{n_1^2}{L_1^2} + \frac{n_2^2}{L_2^2} + \frac{n_3^2}{L_3^2} \right)^{1/2}.
\]
To extract the physically meaningful finite contribution (Casimir energy), 
one employs zeta-function regularization. We define the Epstein zeta function 
for the aspect ratio vector $\mathbf{a} = (1, \sqrt{2}, \sqrt{3})$:
\[
Z(s; \mathbf{a}) = \sum_{\mathbf{n} \in \mathbb{Z}^3 \setminus \{0\}} 
\left( n_1^2 + \frac{n_2^2}{2} + \frac{n_3^2}{3} \right)^{-s}.
\]
The regularized topological coefficient is derived from the analytic continuation 
of $Z(s; \mathbf{a})$ to $s = -1/2$. For irrational tori that satisfy Diophantine 
non-resonance conditions, the analytic continuation is well-behaved and converges 
to a shape-dependent constant \citep{Lutken1986,Kirsten2001}. 

Numerical evaluation of the analytically continued Epstein zeta function for 
these specific irrational ratios yields the topological coefficient 
$C \approx 7.59$. This purely geometric scalar acts as the coupling prefactor 
in the emergent mass relation, eliminating the need for ad-hoc phase corrections. \qed
\end{proof}

%==============================================================================
\section{Numerical Verification}
\label{sec:numeric}

\paragraph{Units.} Throughout this paper we use natural units with $c=\hbar=1$ 
unless otherwise specified. Masses are quoted in $\mathrm{GeV}$, 
understanding $m c^2$ as the corresponding energy.

\begin{table}[h]
\centering
\caption{Parameters for the Higgs mass scaling relation.}
\label{tab:params}
\begin{tabular}{lll}
\toprule
Parameter & Symbol & Value \\
\midrule
Topological scale & $\Ltop$ & $28.57^{+0.73}_{-0.87}~\mathrm{Gpc} = 8.816 \times 10^{26}~\mathrm{m}$ \\
Percolation amplification & $\xiPhi$ & $3.8$ \\
Geometric factor & $\Fgeo$ & $6^{1/6} \approx 1.348$ \\
Mode count & $\Nmodes$ & $5.24 \times 10^{55}$ \\
Topological Casimir coeff. & $C$ & $7.59$ \\
Planck constant & $\hbar$ & $1.0545718 \times 10^{-34}~\mathrm{J\cdot s}$ \\
Gravitational constant & $G$ & $6.67430 \times 10^{-11}~\mathrm{m^3/(kg \cdot s^2)}$ \\
Speed of light & $c$ & $2.99792458 \times 10^8~\mathrm{m/s}$ \\
\bottomrule
\end{tabular}
\end{table}

Substituting into the scaling relation, we split the calculation into numerator $\mathcal{N}$ and denominator $\mathcal{D}$:
\begin{align*}
\mathcal{N} &= 7.59 \cdot 3 \cdot (1.0546 \times 10^{-34}~\mathrm{J\cdot s})^3 \cdot 3.8 \cdot 1.348 \cdot 5.24 \times 10^{55} \\
\mathcal{D} &= 8\pi \cdot 6.674 \times 10^{-11}~\mathrm{m^3/(kg \cdot s^2)} \cdot 3.00 \times 10^8~\mathrm{m/s} \cdot (8.816 \times 10^{26}~\mathrm{m})^2
\end{align*}
Evaluating the ratio yields:
\[
m_H = \left( \frac{\mathcal{N}}{\mathcal{D}} \right)^{1/4} = 125.1~\mathrm{GeV},
\]
in agreement with the PDG value $m_H = 125.10 \pm 0.14~\mathrm{GeV}$.

%==============================================================================
\section{Falsifiable Predictions}
\label{sec:predictions}

The $\ITthree$ scaling ansatz makes sharp, testable predictions:

\begin{enumerate}
    \item \textbf{Redshift evolution of $m_H$:} Since $\xiPhi$ depends on the void 
    fraction $\phi(z)$, the effective Higgs mass should exhibit weak evolution:
    \[
    m_H(z) \approx m_H(0) \cdot \left[1 + \epsilon \cdot e^{-(z/0.5)^2}\right], 
    \quad \epsilon \sim 10^{-3}.
    \]
    Future high-precision measurements at FCC-ee could detect this signature.

    \item \textbf{Correlation with CMB anomalies:} Fluctuations in the low-$\ell$ 
    CMB power spectrum should correlate with inferred variations in $m_H$ at 
    recombination, testable with CMB-S4 data.

    \item \textbf{Universality across masses:} Analogous formulas should derive 
    all Standard Model masses from $\Ltop$. For example, the electron mass:
    \[
    m_e = \left\{ \frac{C_e \cdot 3\hbar^3 \, \xiPhi \, \Fgeo \, \Nmodes^{(e)}}{8\pi G \, c \, \Ltop^2} \right\}^{1/4},
    \]
    with mode count $\Nmodes^{(e)}$ determined by the electron's Compton wavelength. 
    Verification across multiple particles would confirm the paradigm.

    \item \textbf{Topological corrections at high precision:} The geometric 
    constant $C$ may receive small corrections from higher-order topology. 
    Measurements of $m_H$ at the $10^{-4}$ level (FCC-ee target) could reveal 
    these effects.

    \item \textbf{Cross-scale percolation signature:} The same percolation 
    mathematics that enters the mass scaling governs quark-gluon plasma formation 
    at the LHC \citep{Logvinovich2026i}. Correlations between heavy-ion collision 
    observables and cosmological parameters would provide independent validation.
\end{enumerate}

Failure to detect these signatures at the predicted levels would falsify the 
$\ITthree$ derivation.

%==============================================================================
\section{Conclusion}
\label{sec:conclusion}

We have presented a phenomenological scaling relation that connects the Higgs 
boson mass to the topological scale of the universe within the $\ITthree$ paradigm. 
The derivation rests on three independently motivated pillars: holographic vacuum 
energy, percolation spectral dimension ($d_s = 4/3$), and topological Casimir 
regularization via the Epstein zeta function on irrational tori. 

We explicitly frame this as an \textit{emergent effective scaling law}, not a 
first-principles quantum field theory derivation. Its strength lies in predictive 
economy: one cosmological parameter $\Ltop$, zero new particles, zero modifications 
to General Relativity, and sharp falsifiable predictions for $m_H(z)$ evolution, 
tensor modes, and cross-scale percolation signatures.

If future data from FCC-ee, CMB-S4, and Euclid confirm these predictions, the 
hierarchy problem will be recast as a geometric consequence of cosmic finitude. 
If not, the ansatz will be refined or discarded. That is the standard of 
empirical science.

The universe may not require invisible substances to explain its scales. It may 
only require finite geometry, critical connectivity, and the mathematics of emergence.

%==============================================================================
\section*{Outlook: From Scaling Ansatz to Fundamental Theory}

The derivation presented here should be viewed not as a final explanation, but as 
a \textit{phenomenological bridge} between cosmic topology and particle physics. 
Its success would motivate three concrete theoretical directions:

\begin{enumerate}
    \item \textbf{Effective field theory on compact manifolds}: Derive the 
    Higgs potential $V(\phi_H)$ in a toroidal geometry with topological boundary 
    conditions, showing how vacuum expectation value $v$ emerges from $\Ltop$.
    
    \item \textbf{Spectral geometry of the irrational torus}: Compute the 
    regularized Casimir energy for the Standard Model field content on 
    $\mathbb{T}^3(1,\sqrt{2},\sqrt{3})$, verifying whether the coefficient 
    $C \approx 7.59$ arises from first principles.
    
    \item \textbf{Percolation-renormalization group flow}: Establish whether 
    the spectral dimension $d_s = 4/3$ of critical percolation clusters 
    corresponds to a fixed point in the RG flow of an effective scalar theory 
    coupled to topology.
\end{enumerate}

Conversely, if precision measurements falsify the scaling predictions, the 
ansatz will require revision---perhaps pointing to additional topological 
invariants or a different percolation universality class. 

Either outcome advances understanding. That is the power of a well-posed 
phenomenological hypothesis.

%==============================================================================
\section*{Acknowledgments}
We thank the Planck, ATLAS, CMS, and ALICE collaborations for open data access. 
This research utilized CLASS, NumPy, SciPy, and Matplotlib. Computations were 
performed on a local Apple Intel-based MacBook Pro workstation. The author 
acknowledges the pioneering work on cosmic topology by Jean-Pierre Luminet, 
on percolation theory by Dietrich Stauffer, and on the hierarchy problem by 
Leonard Susskind and Steven Weinberg.

%==============================================================================
\begin{thebibliography}{99}

\bibitem[Alexander \& Orbach(1982)]{Alexander1982}
Alexander, S., \& Orbach, R. (1982). \textit{Density of states on fractals: fractons}. Journal de Physique Lettres, 43(17), L625-L631.

\bibitem[Kirsten(2001)]{Kirsten2001}
Kirsten, K. (2001). \textit{Spectral Functions in Mathematics and Physics}. Chapman and Hall/CRC.

\bibitem[Logvinovich(2026a)]{Logvinovich2026a}
Logvinovich, V. (2026a). \textit{Flat Irrational Torus Topology: Simultaneous Resolution of Low-$\ell$ Anomaly and Hubble Tension} (Paper I). Zenodo. DOI: 10.5281/zenodo.19431323.

\bibitem[Logvinovich(2026b)]{Logvinovich2026b}
Logvinovich, V. (2026b). \textit{Topological Energy Sink and Thermodynamic Stability} (Paper II). Zenodo. DOI: 10.5281/zenodo.19473353.

\bibitem[Logvinovich(2026c)]{Logvinovich2026c}
Logvinovich, V. (2026c). \textit{The Cosmic Sponge: Percolation Amplification in Large-Scale Structure} (Paper III). Zenodo. DOI: 10.5281/zenodo.19479551.

\bibitem[Logvinovich(2026d)]{Logvinovich2026d}
Logvinovich, V. (2026d). \textit{The Cosmic Sponge Mechanism: Late-Time Activation Resolves the Hubble Tension} (Paper IV). Zenodo. DOI: 10.5281/zenodo.19489307.

\bibitem[Logvinovich(2026e)]{Logvinovich2026e}
Logvinovich, V. (2026e). \textit{Topological Damping of Structure Growth: Resolving the S8 Tension} (Paper V). Zenodo. DOI: 10.5281/zenodo.19489122.

\bibitem[Logvinovich(2026f)]{Logvinovich2026f}
Logvinovich, V. (2026f). \textit{Geometric Acceleration Scale and Flat Rotation Curves} (Paper VI). Zenodo. DOI: 10.5281/zenodo.19489307.

\bibitem[Logvinovich(2026g)]{Logvinovich2026g}
Logvinovich, V. (2026g). \textit{Holographic Vacuum Energy and the $10^{120}$ Discrepancy} (Paper VII). Zenodo. DOI: 10.5281/zenodo.19489438.

\bibitem[Logvinovich(2026h)]{Logvinovich2026h}
Logvinovich, V. (2026h). \textit{Topological Genesis of Primordial Perturbations: Inflation without Inflaton} (Paper VIII). Zenodo. DOI: 10.5281/zenodo.19489594.

\bibitem[Logvinovich(2026i)]{Logvinovich2026i}
Logvinovich, V. (2026i). \textit{Universal Percolation and Collective Flow: Connecting QGP to the Cosmic Sponge} (Paper IX). Zenodo. DOI: 10.5281/zenodo.19492202.

\bibitem[Logvinovich(2026j)]{Logvinovich2026j}
Logvinovich, V. (2026j). \textit{The IT3 Paradigm: A Unified Geometric Framework for Cosmology and Fundamental Physics} (Paper X). Zenodo. DOI: 10.5281/zenodo.19492203.

\bibitem[L\"utken \& Ordo\~nez(1986)]{Lutken1986}
L\"utken, C. A., \& Ordo\~nez, C. R. (1986). \textit{Vacuum energy of Kaluza-Klein spaces}. Journal of Physics A: Mathematical and General, 19(17), 3505.

\bibitem[Stauffer \& Aharony(1994)]{Stauffer1994}
Stauffer, D., \& Aharony, A. (1994). \textit{Introduction to Percolation Theory} (2nd ed.). Taylor \& Francis.

\end{thebibliography}

%==============================================================================
\appendix
\section{Reproducibility and Code Availability}
\label{app:reproducibility}

All code, MCMC chains, and mathematical proofs supporting this paper are 
available in the public repository:
\begin{center}
\url{https://github.com/Viktar-Pi/FlatIrrationalTorus}
\end{center}
The repository includes:
\begin{itemize}
    \item \texttt{analysis/} --- MCMC scripts (emcee + CLASS wrapper)
    \item \texttt{data/} --- Planck PR4 likelihood files and derived chains
    \item \texttt{figures/} --- Generated plots (corner, correlation, validation)
    \item \texttt{proofs/} --- Mathematical derivations in PDF format
    \item \texttt{Dockerfile} --- Reproducible containerized environment
    \item \texttt{README.md} --- Step-by-step execution guide
\end{itemize}
Running \texttt{python3 run\_all.py} reproduces all results and figures 
reported in this manuscript. The permanent DOI for this work is: 
\texttt{10.5281/zenodo.XXXXXXX}.

\end{document}